% HD-I18. GPT-6 Astra Pro, /root/imprimitive, inherited effort.
% 21 November 2026. Exact remaining nonreflection stratum.
\subsection{HD-I18: two fully colored branch fibres give an exact immersion}
\label{sec:HD-I18}

Retain the actual tower of HD-I01,
\[
 X=Z/\ker\chi\longrightarrow Y=Z/H
                 \xrightarrow{h}\mathbf P^1,
 \qquad \deg h=m,
\]
for the natural representation of $G(d,e,m)$, with
$d\in\{3,4,6\}$, $e\mid d$, balanced signature $(n,n)$ and
$n\geq5$. A permutation cycle of a monomial matrix is called
colored when the product of its nonzero entries is not one.
The characteristic polynomial of a cycle of length $a$ and color
$s\in\mathbf Z/d$ is $T^a-\zeta_d^s$. Consequently all its
cycles are colored exactly when the inertia matrix has no fixed
vector in the natural representation.

Consider the branch profiles having exactly two nonreflection
values, every cycle above either of them colored, and all remaining
branch matrices reflections. On every Nielsen component of these
profiles the forgetful map to the cyclic-cover data $(X\to Y)$ has
injective differential at every point. Its rank is exactly
\begin{equation}
                         2n-1.                       \tag{I18.1}
\end{equation}
For the components with $g(Y)\leq n-2$, the Weil period map also
has rank exactly $2n-1$ at every point. All statements refer to
local marked Hurwitz spaces, or equivalently to a finite level
cover, so finite automorphisms introduce no tangent ambiguity.

These profiles occur for every $d,e,n,m$ in the prompt. More
precisely, the explicit tuple in paragraph~3 below generates the
full group $G(d,e,m)$, has $g(Y)=n-2$, has six visible cyclic
branch points, and has period rank $2n-1$ everywhere on each
Nielsen component containing it. In particular, for $n=5$ this
gives exact rank-nine components for every allowed imprimitive
group and every allowed $m$.

\paragraph{1. The two complete fibres and their local constraints.}
Let $b_1,b_2$ be the two nonreflection values, and put
$E=h^*T_{\mathbf P^1}$. A fixed-vector-free monodromy has
codimension of fixed vectors equal to $m$; every reflection has
codimension one. The balanced parabolic dimension formula of
HD-I05 therefore makes the number of reflection values exactly
$2n$. The total branch count is $b=2n+2$, so the Hurwitz
dimension is $b-3=2n-1$.

Use the exact jet map $J$ of HD-I16. At a colored cycle of length
$a$, the marked point is fixed while the branch value moves.
In a local coordinate with $h(z)=b_i+z^a$, a variation
$q(z)=q_0+q_1z+\cdots$ must satisfy
\[
 q_1=\cdots=q_{a-1}=0,
\]
and the constant terms at all points over the same $b_i$ must
agree in $T_{b_i}\mathbf P^1$. There are no uncolored points in
these two fibres. Thus the jet divisor is exactly
\begin{equation}
 D_{\rm jet}=h^*(b_1+b_2),\qquad
 E(-D_{\rm jet})\simeq\mathcal O_Y.                  \tag{I18.2}
\end{equation}
It has length $2m$, while the constraint space has dimension
$N=2m-2$. Its two-dimensional kernel $U$ consists of one
constant target tangent value on each complete fibre:
\[
 0\longrightarrow U\longrightarrow H^0(E|_{D_{\rm jet}})
                  \xrightarrow{\pi}\mathbf C^{2m-2}
                  \longrightarrow0.
\]
The map in HD-I16 is $J=\pi\circ\operatorname{ev}$.

Evaluation of global vector fields on $\mathbf P^1$ at the two
distinct points is surjective. Pulling such a vector field back
by $h$ gives locally $v(b_i+z^a)=v(b_i)\bmod z^a$.
Consequently every vector of $U$ is the restriction of a
pulled-back global target vector field. On the other hand,
\[
 \ker\operatorname{ev}
   =H^0(E(-D_{\rm jet}))=H^0(\mathcal O_Y)
\]
has dimension one. It is exactly the pullback of the line of
target vector fields vanishing at both $b_1,b_2$. This proves
the equality of actual spaces
\begin{equation}
       \ker J=h^*H^0(T_{\mathbf P^1}),\qquad
                         \dim\ker J=3.              \tag{I18.3}
\end{equation}
In particular the kernel is exhausted by target coordinate
changes, with no residual deformation of the map at fixed
cyclic data. This proof of (I18.3) does not require a
jet-surjectivity estimate.

The cyclic datum is infinitesimally the marked curve $(Y,R)$:
the root line bundle and its prescribed cyclic powers are finite
etale data in characteristic zero. The stability inequality is
$2g-2+|R|=2n>0$, so it has no infinitesimal automorphisms.
The exact tangent-fibre description of HD-I16 quotients (I18.3)
by precisely $H^0(T_{\mathbf P^1})$. Hence the forgetful
differential is injective, proving (I18.1).

For comparison with the numerical jet formula, the prompt's
bound on $m$ gives $H^1(E)=0$. Therefore
\begin{equation}
 \operatorname{rank}J=h^0(E)-3=2m-g-2,
 \qquad N-\operatorname{rank}J=g.                    \tag{I18.4}
\end{equation}
Thus every positive-genus component here has a non-surjective
constraint map. The criterion $\deg D_{\rm jet}\leq2m-2g+1$
from HD-I16 does not apply for $g\geq1$, whereas (I18.2)--(I18.4)
compute its exact deficiency and forgetful rank.

\paragraph{2. Exact period rank in the indicated genera.}
For $g\leq n-2$, HD-I12 proves at every cyclic datum the
surjectivity of the actual multiplication map
\[
 H^0(A)\otimes H^0(B)
       \longrightarrow H^0(\omega_Y^2(R)),\qquad
 A=\omega_Y\otimes L_\chi,
 \quad B=\omega_Y\otimes L_{\bar\chi}.
\]
Its dual is the cyclic Weil-period differential, up to the
fixed nonzero polarization scalar identified in HD-I01 and
HD-I05. It is therefore injective. Composing it with the
injective forgetful differential proves the period assertion,
including its all-point, individual-component scope. No claim
about the multiplication kernel in the remaining genera is
needed here.

\paragraph{3. Explicit nonempty full-monodromy components.}
Put $a=m-2$, and let $A$ act on the standard basis by
\[
 e_j\longmapsto e_{j+1}\ (1\leq j<a),\qquad
 e_a\longmapsto\zeta_d e_1,\qquad
 e_{m-1}\longmapsto\zeta_d^u e_{m-1},\qquad
 e_m\longmapsto\zeta_d^v e_m.
\]
Choose the nonzero colors from the following complete table:
\[
\begin{array}{c|c|c|c}
 d&e&u&v\\\hline
 3&1&1&2\\
 3&3&1&1\\
 4&1&1&3\\
 4&2&2&3\\
 4&4&1&2\\
 6&1&1&5\\
 6&2&2&5\\
 6&3&1&1\\
 6&6&1&4
\end{array}
\]
In each row $1+u+v\equiv e\pmod d$, with $e=d$
interpreted as zero. Thus $A\in G(d,e,m)$, and its three
cycles have nonzero colors $1,u,v$. The three inverse cycles
of $A^{-1}$ have the complementary nonzero colors.

Let $s=(1\,2)$, $t=(1\,m-1)$ and $w=(1\,m)$ be
uncolored transposition matrices. Let $s'$ be the colored
transposition
\[
 e_1\longmapsto\zeta_d e_2,\qquad
 e_2\longmapsto\zeta_d^{-1}e_1,
\]
fixing all other basis vectors. All four matrices are involutory
reflections in $G(d,e,m)$. Take the ordered tuple
\begin{equation}
 \bigl(A,A^{-1},s,s,t,t,w,w,s',s',
                 \underbrace{s,s,\ldots,s,s}_{(n-4)\ \text{pairs}}\bigr).
                                                        \tag{I18.5}
\end{equation}
It has product one and contains exactly $2n$ reflection entries.
The prompt gives $m\geq9$, so the displayed indices and the
long cycle are distinct as required.

Here is a full group-generation proof. For $0\leq j\leq a-2$,
\[
             A^jsA^{-j}=(j+1\,j+2)
\]
is uncolored: both basis vectors have moved along the long cycle
without passing its colored closing edge. These adjacent
transpositions generate the full uncolored symmetric group on
$\{1,\ldots,a\}$. Adding $t,w$ gives the full uncolored
$\mathfrak S_m$. The diagonal matrix $s's$ has colors
$(-1,1,0,\ldots,0)$. Its symmetric-group conjugates generate
the entire diagonal augmentation subgroup
\[
 \{(c_1,\ldots,c_m)\in(\mathbf Z/d)^m:
                                      \textstyle\sum c_i=0\}.
\]
Indeed every such vector is a sum of $c_j(e_j-e_1)$,
$2\leq j\leq m$. Modulo this subgroup and the symmetric
group, total color identifies the remaining quotient with
$\mathbf Z/d$. The total color of $A$ generates exactly
$e\mathbf Z/d$. Thus the generated group is all of
$G(d,e,m)$, including the cases $e=d$.

The permutation action is transitive, so its degree-$m$
intermediate curve is connected. The two long cycles have
ramification deficits $m-3$ each; each reflection contributes
one. Riemann--Hurwitz gives
\[
 2g(Y)-2=-2m+2(m-3)+2n=2n-6,
 \qquad g(Y)=n-2.
\]
There are exactly six visible cyclic branch points, from the
three colored cycles of $A$ and of $A^{-1}$. Their fractional
colors sum to three, since they occur in complementary pairs.
The actual line-stabilizer character has order $d$: the
augmentation subgroup already contains a diagonal matrix with
first entry $\zeta_d$. The cyclic cover is therefore connected,
and its two primitive Hodge dimensions are
\[
                    g(Y)-1+3=n.
\]
This verifies the balanced signature through the exact tower,
not just through a dimension count.

Finally, the product-one generating tuple gives a connected
branched cover for every configuration of distinct branch
values: glue copies of a cut sphere using its sheet permutations,
and use the local charts $z\mapsto z^{\operatorname{ord}(\gamma)}$
to compactify. The resulting compact Riemann surface is projective.
Thus (I18.5) supplies actual nonempty Hurwitz components for
every displayed group, not merely formal branch data.

\paragraph{Checks and scope.}
The author read the entire proof, including the exact complete-fibre
jet sequence, the actual target-vector-field kernel, the period
composition and the full colored group-generation argument.
The separate exact finite certificate
\texttt{two\_full\_fiber\_certificate.py} checks the literal
monomial identities, cycle colors, signature and genus for
54 parameter choices. The universal group and jet statements
are proved above; they are not inferred from these samples.
The exact rank belongs to the stated individual nonreflection
components. Other branch profiles retain their own multiplication
and jet maps from HD-I16.
